\documentclass[12pt]{article}
\usepackage[utf8]{inputenc}
\usepackage[T1]{fontenc}
\usepackage[french]{babel}
\usepackage{amsmath, amssymb, amsthm}
\usepackage{hyperref}
\usepackage{geometry}
\geometry{a4paper, margin=1in}

\title{Analyse et R\'esultats Partiels sur la Conjecture d'Erd\H{o}s sur les Progressions Arithm\'etiques}
\author{Charles EDOU NZE\thanks{Charles EDOU NZE, chercheur ind\'ependant}}
\date{}

\newtheorem{theorem}{Th\'eor\`eme}
\newtheorem{lemma}[theorem]{Lemme}
\newtheorem{definition}[theorem]{D\'efinition}

\begin{document}

\maketitle

\begin{abstract}
Ce document fournit une d\'ecomposition axiomatique rigoureuse, une revue de la litt\'erature et une structure de preuve partielle \'etape par \'etape pour la Conjecture d'Erd\H{o}s sur les Progressions Arithm\'etiques. L'architecture est con\c{c}ue pour faciliter une future autoformalisation dans des syst\`emes tels que Lean 4.
\end{abstract}

\section{Analyse et D\'ecomposition}

\subsection{D\'efinitions Axiomatiques}
Soit $\mathbb{N} = \{1, 2, \dots \}$ l'ensemble des entiers strictement positifs. Nous d\'efinissons les types et les variables de mani\`ere rigoureuse.
\begin{definition}[Somme Harmonique Divergente]
Soit $A \subseteq \mathbb{N}$ un sous-ensemble. Nous disons que $A$ poss\`ede une somme harmonique divergente, not\'ee $\mathcal{D}(A)$, si la s\'erie des inverses de ses \'el\'ements diverge :
\[
\sum_{n \in A} \frac{1}{n} = \infty.
\]
\end{definition}

\begin{definition}[Progression Arithm\'etique de Longueur $k$]
Pour $k \in \mathbb{N}, k \ge 3$, un ensemble $A \subseteq \mathbb{N}$ contient une progression arithm\'etique de longueur $k$, not\'ee $\mathcal{AP}_k(A)$, s'il existe $a \in \mathbb{N}$ et $d \in \mathbb{N}$ tels que pour tout $0 \le i < k$, $a + id \in A$.
\end{definition}

\textbf{\'Enonc\'e de la Conjecture :}
Pour tout $A \subseteq \mathbb{N}$, si $\mathcal{D}(A)$ est vraie, alors pour tout $k \ge 3$, $\mathcal{AP}_k(A)$ est vraie.

\subsection{Structures Alg\'ebriques et Combinatoires Sous-jacentes}
Le probl\`eme m\^ele fondamentalement combinatoire additive et th\'eorie analytique des nombres. L'ensemble $A$ est consid\'er\'e comme un sous-ensemble dense au sens pond\'er\'e. La structure naturelle implique l'analyse de Fourier sur les groupes cycliques $\mathbb{Z}/N\mathbb{Z}$ (m\'ethode du cercle de Hardy-Littlewood) et l'utilisation des normes de Gowers pour d\'etecter les structures lin\'eaires. Une perspective alternative utilise la th\'eorie ergodique, associant le sous-ensemble $A$ \`a un syst\`eme dynamique pr\'eservant la mesure.

\section{Recherche de Litt\'erature Contextuelle}
Le th\'eor\`eme le plus important li\'e \`a cette conjecture est le th\'eor\`eme de Szemer\'edi, qui \'etablit que tout sous-ensemble de $\mathbb{N}$ ayant une densit\'e sup\'erieure strictement positive contient des progressions arithm\'etiques de longueur arbitraire. Le th\'eor\`eme de Green-Tao constitue une avanc\'ee majeure en prouvant que les nombres premiers, un ensemble de densit\'e nulle mais de somme harmonique divergente, contiennent des progressions arithm\'etiques de longueur arbitraire.
Les r\'ecentes avanc\'ees de Kelley et Meka (2023) sur le th\'eor\`eme de Roth donnent des bornes quantitatives fortes sur les ensembles sans progression arithm\'etique de longueur 3.

Par analogie, la r\'esolution du probl\`eme du Cap Set par Ellenberg et Gijswijt utilise la m\'ethode polynomiale. Bien que les m\'ethodes des corps finis ne se traduisent pas directement dans $\mathbb{Z}$, les m\'ethodes alg\'ebriques analogues aux bornes de Croot-Lev-Pach fournissent des indications pour borner les sous-ensembles sans structures lin\'eaires.

\section{Strat\'egie de Preuve et Isolation des Lemmes}
Pour aborder une version simplifi\'ee (par exemple, $k=3$ sur des sous-ensembles restreints), nous d\'ecomposons le probl\`eme en plusieurs lemmes.

\begin{itemize}
    \item \textbf{Lemme 1 (Incr\'ement de Densit\'e) :} Si un sous-ensemble $A$ de $\{1, \dots, N\}$ ne contient pas de progression arithm\'etique de longueur 3, il doit \^etre corr\'el\'e avec un ensemble de Bohr, ce qui permet un incr\'ement de densit\'e sur une sous-structure structur\'ee.
    \item \textbf{Lemme 2 (Borne Harmonique sur la Clairsemance) :} Si un ensemble $A$ ne contient aucune progression arithm\'etique de longueur 3, sa fonction de comptage $A(x) = |A \cap \{1, \dots, x\}|$ est born\'ee par $O(x / (\log x)^c)$ pour un certain $c > 1$.
\end{itemize}

\section{Preuve Informelle du Lemme 2}
\begin{proof}[Preuve du Lemme 2]
Nous pr\'esentons la preuve \'etape par \'etape. Soit $A \subset \mathbb{N}$ tel que $A$ ne contient aucune progression arithm\'etique de longueur 3.
Soit $N$ un grand entier, et soit $A_N = A \cap \{1, 2, \dots, N\}$.
D'apr\`es les bornes quantitatives de Kelley et Meka, la taille de $A_N$ v\'erifie l'in\'egalit\'e :
\[
|A_N| \le \frac{N}{\exp(c (\log N)^{\beta})}
\]
pour certaines constantes $c > 0$ et $\beta > 0$.
Nous \'evaluons la somme harmonique de $A$ par sommation d'Abel.
Soit $S(N) = \sum_{n \in A, n \le N} \frac{1}{n}$.
En utilisant la sommation par parties, nous exprimons $S(N)$ en fonction de $A(x) = |A \cap \{1, \dots, \lfloor x \rfloor\}|$ :
\[
S(N) = \frac{A(N)}{N} + \int_{1}^{N} \frac{A(t)}{t^2} dt.
\]
Nous majorons le premier terme :
\[
\frac{A(N)}{N} \le \frac{1}{\exp(c (\log N)^{\beta})}.
\]
Lorsque $N$ tend vers l'infini, ce terme tend vers $0$.
Pour le terme int\'egral, nous substituons la majoration de $A(t)$ :
\[
\int_{1}^{N} \frac{A(t)}{t^2} dt \le \int_{1}^{N} \frac{t \exp(-c (\log t)^{\beta})}{t^2} dt = \int_{1}^{N} \frac{1}{t \exp(c (\log t)^{\beta})} dt.
\]
Nous calculons cette int\'egrale \`a l'aide du changement de variable $u = \log t$. Ainsi, $du = \frac{1}{t} dt$. Les bornes d'int\'egration changent de $t=1$ \`a $u=0$, et de $t=N$ \`a $u = \log N$ :
\[
\int_{0}^{\log N} \frac{1}{\exp(c u^{\beta})} du.
\]
Pour tout $\beta > 0$, l'int\'egrale $\int_{0}^{\infty} \exp(-c u^{\beta}) du$ est convergente. Plus pr\'ecis\'ement, en posant $v = c u^{\beta}$, $dv = c \beta u^{\beta - 1} du$, nous trouvons que l'int\'egrale est proportionnelle \`a la fonction Gamma $\Gamma(1/\beta)$.
Puisque l'int\'egrale est major\'ee par une constante finie ind\'ependante de $N$, il s'ensuit que $S(N)$ est born\'e lorsque $N \to \infty$.
Ainsi, $\sum_{n \in A} \frac{1}{n} < \infty$.
Par contrapos\'ee, si $\sum_{n \in A} \frac{1}{n} = \infty$, l'ensemble $A$ doit contenir une progression arithm\'etique de longueur 3. Ceci ach\`eve la preuve du Lemme 2.
\end{proof}

\section{Architecture pour l'Autoformalisation}
La structure de la preuve est directement adapt\'ee pour une impl\'ementation dans Lean 4.

\begin{verbatim}
import Mathlib.Data.Real.Basic
import Mathlib.Analysis.SpecialFunctions.Log.Basic

-- Axiomatic definition of the divergent harmonic sum
def has_divergent_harmonic_sum (A : Set Nat) : Prop :=
  Filter.Tendsto (fun N => \sum_{n \in A \cap {x | x \le N}} (1 / (n : Real)))
    Filter.atTop Filter.atTop

-- Axiomatic definition of containing an arithmetic progression of length k
def contains_AP (A : Set Nat) (k : Nat) : Prop :=
  \exists a d : Nat, d > 0 \and \forall i : Nat, i < k \to a + i * d \in A

-- Main Theorem Statement (Erdős Conjecture)
theorem erdos_ap_conjecture (A : Set Nat) (h : has_divergent_harmonic_sum A) :
  \forall k \ge 3, contains_AP A k := by
  sorry

-- Lemma 2
lemma ap3_harmonic_bound (A : Set Nat) (h : \not contains_AP A 3) :
  \not has_divergent_harmonic_sum A := by
  sorry
\end{verbatim}

\end{document}
